\documentclass[11pt,letterpaper]{article}

% --- Packages ---
\usepackage[top=1in,bottom=1in,left=1in,right=1in]{geometry}
\usepackage{amsmath,amssymb}
\usepackage{booktabs}
\usepackage{array}
\usepackage{xcolor}
\usepackage{enumitem}
\usepackage{hyperref}
\usepackage{microtype}
\documentclass{article}
\usepackage{graphicx} % Required to handle images
% --- Hyperref Setup ---
\hypersetup{
    colorlinks=true,
    linkcolor=blue!70!black,
    citecolor=blue!70!black,
    urlcolor=blue!70!black
}

% --- Styling Commands ---
\newcommand{\openitem}[1]{\textbf{\textcolor{red!75!black}{[OPEN: #1]}}}
\setlist[enumerate]{leftmargin=1.8em, itemsep=2pt, topsep=3pt}
\setlist[itemize]{leftmargin=1.5em, itemsep=2pt, topsep=3pt}

\title{\textbf{METHODS DRAFT: Algorithmic Calibration and Micro-Randomized Decision Engine for Just-in-Time Adaptive Interventions: Methods}}
\author{Tien (Tyler) Le}
\date{September 2026}

\begin{document}

\maketitle

\section{MSSD Definition and Decision Rules}

\subsection{Mathematical Definition}

Let $x_{i,t} \in \{1, 2, \dots, 7\}$ denote participant $i$'s integer Likert response to a target Ecological Momentary Assessment (EMA) item at the $t$-th completed prompt. The full-series mean square of successive differences (MSSD) is conventionally defined as:
\begin{equation}
\text{MSSD}_i = \frac{1}{n_i - 1} \sum_{t=2}^{n_i} (x_{i,t} - x_{i,t-1})^2
\end{equation}
In real-time deployment, evaluating a full-series metric is contraindicated because it averages over historical baseline shifts rather than isolating acute proximal swings. The decision engine evaluates a rolling, recency-weighted formulation over a trailing window of $w = 3$ successive differences:
\begin{equation}
\widehat{\text{MSSD}}_{i,t} = \frac{1}{\min(w,\, t - 1)} \sum_{k=\max(2,\, t - w + 1)}^{t} (x_{i,k} - x_{i,k-1})^2
\end{equation}
with $w = 3$ and a minimum evaluation period of 1. At $t = 2$ and $t = 3$, the engine computes the metric over available differences rather than remaining undefined. Offline analytical audits must reference the runtime telemetry logged in \texttt{JITAILog.observed\_mssd} rather than recomputing a global retrospective MSSD, as the rolling and static formulations yield systematically distinct quantities.

\subsection{Production Decision Pipeline}

At each completed EMA prompt, the decision engine executes a sequential gating cascade:
\begin{enumerate}
    \item \textbf{Burn-in Initialization} 
    
    The initial observation ($t = 1$, where t is the number of days) provides no prior difference. Prompts $t = 2, 3, 4$ yield rolling MSSD estimates but do not evaluate trigger eligibility, as the within-person threshold requires a minimum of three historical MSSD observations. Randomized decisions become active strictly from prompt $t \ge 5$ or week 2 onward.
    
    \item \textbf{Expanding Within-Person Threshold0} 
    
    The intervention trigger cutoff is calibrated to the 80th percentile of participant $i$'s expanding $\widehat{\text{MSSD}}$ vector:
    \begin{equation}
    \Theta_{i,t} = \text{Quantile}_{0.80}\left(\{\widehat{\text{MSSD}}_{i,1}, \dots, \widehat{\text{MSSD}}_{i,t-1}\}\right)
    \end{equation}
    Lagging the evaluation window by one prompt ensures strictly out-of-sample benchmarking, preventing the current affective deviation from inflating its own threshold.
    
    \item \textbf{Eligibility Gate} 
    
    An intervention decision point is flagged as eligible if and only if:
    \begin{equation}
    \widehat{\text{MSSD}}_{i,t} > \Theta_{i,t}
    \end{equation}
    Interventions are triggered by state volatility relative to an idiographic baseline rather than absolute distress levels.
    
    \item \textbf{Daily Dosage Cap} 
    
    A strict ceiling limits deliveries to no more than 4 prompts per calendar day. Decision points evaluated after reaching the daily cap are logged as blocked.
    
    \item \textbf{Refractory Cooldown} 
    
    A minimum interval of *30 minutes* must elapse following a delivered prompt. Subsequent threshold exceedances within this refractory window are suppressed to prevent notification habituation.
    
    \item \textbf{Micro-Randomization} 
    
    Conditioned on clearing the eligibility, daily cap, and cooldown gates, intervention delivery is randomized via a Bernoulli trial with fixed probability $p = 0.50$:
    \begin{equation}
    A_{i,t} \sim \text{Bernoulli}(0.50)
    \end{equation}
    Delivery randomization uncouples causal intervention effects from the underlying deterministic gating mechanism.
\end{enumerate}

\section{Synthetic Data Generation}

\subsection{Generative Process}

Synthetic cohorts are simulated using a stationary, discrete-time, first-order autoregressive ($\text{AR}(1)$) Gaussian process parameterized independently per participant $i$ and construct $c$:
\begin{align}
z_{i,1} &\sim \mathcal{N}(0,\, \sigma_i^2) \\
z_{i,t} &= \rho_i z_{i,t-1} + \varepsilon_{i,t}, \quad \varepsilon_{i,t} \sim \mathcal{N}\left(0,\, \sigma_i^2(1 - \rho_i^2)\right) \\
x_{i,t} &= \text{clip}\left(\left\lfloor \mu_i + z_{i,t} \right\rceil,\, 1,\, 7\right)
\end{align}
Scaling the innovation variance by $\sigma_i^2(1 - \rho_i^2)$ preserves a constant marginal variance $\text{Var}(z_{i,t}) = \sigma_i^2$ for all $t$. The latent process reverts to set point $\mu_i$ without autonomous deterministic drift. The synthetic EMA, $x_i, t$ is clipped to a Likert scale of 1-7 to simulate an EMA response.

\subsection{Idiographic Parameter Sampling}

Participant-level parameters are sampled from continuous uniform distributions across the production affective items (Energy = \texttt{B1}, Stress = \texttt{B2}):

\begin{table}[htbp]
\centering
\small
\begin{tabular}{@{}lccc@{}}
\toprule
\textbf{Construct} & \textbf{Baseline Set Point ($\mu$)} & \textbf{Stationary Volatility ($\sigma$)} & \textbf{Autoregressive Inertia ($\rho$)} \\
\midrule
Energy (\texttt{B1}) & $\mathcal{U}(3.0,\, 5.0)$ & $\mathcal{U}(0.4,\, 1.4)$ & $\mathcal{U}(0.2,\, 0.8)$ \\
Stress (\texttt{B2}) & $\mathcal{U}(2.0,\, 4.0)$ & $\mathcal{U}(0.4,\, 1.2)$ & $\mathcal{U}(0.2,\, 0.8)$ \\
\bottomrule
\end{tabular}
\caption{Continuous parameter bounds for synthetic participant generation.}
\end{table}

\subsection{Clustered Missing Data Architecture}

Missingness is generated via a two-state Markov chain to reflect bursty, dependent non-response patterns rather than an independent Missing Completely at Random (MCAR) mechanism. To put simply, EMA response missingness has a 'sticky' factor. Consecutive uncompleted assessments form clustered sequences with an empirical mean gap length of 3 prompts ($\lambda_{\text{gap}} = 3$) calibrated to a target global response rate of $R = 0.80$.

\subsection{Scenario Taxonomy}

The core simulation pipeline utilizes continuous parameter sampling rather than discrete participant clusters. Controlled scenarios enter through two downstream protocols:
\begin{itemize}
    \item \textbf{Analytical Volatility Stratification} 
    
    For descriptive benchmarking, participants are partitioned into Stable ($E[\text{MSSD}] \le 0.4$), Moderate ($0.4 < E[\text{MSSD}] < 1.6$), and Volatile ($E[\text{MSSD}] \ge 1.6$) strata using the analytical expectation:
    \begin{equation}
    E[\text{MSSD}_i] = 2\sigma_i^2(1 - \rho_i)
    \end{equation}
    \item \textbf{Dynamic Trajectory Scenarios0} 
    
    In misspecification testing, stationary standard deviations $\sigma$ are replaced by dynamic variance schedules $\sigma_t$ to evaluate non-stationary regime adaptation.
\end{itemize}

\subsection{Intraday Delivery Architecture}

The canonical study duration yields $N = D \times M$ scheduled prompts, where $D = 35$ days and $M = 5$ prompts per day. Prompts are distributed across a uniform intraday grid (09:00, 11:00, 13:00, 15:00, 17:00). Each notification carries a 60-minute response expiration window; completed prompts follow an empirical latency distribution with 88\% submitted within 3--55 minutes and 12\% submitted in a delayed window (61--180 minutes).

\subsection{Analytical Ground Truth and Schema Specification}

Ground truth for statistical validation is established analytically via the continuous stationary expectation of the underlying $\text{AR}(1)$ generative process:
\begin{equation}
E[\text{MSSD}_i] = 2\sigma_i^2(1 - \rho_i)
\end{equation}
Algorithmic parameter and threshold recovery are evaluated directly against this latent closed-form benchmark, eliminating reliance on subjective event thresholds or heuristic crisis labels.

To ensure computational reproducibility and seamless translation to production infrastructure, the simulation architecture cross-references two complementary schema specifications:
\begin{enumerate}
    \item \textbf{Conceptual Data Dictionary (\texttt{syntheticData/SCHEMA.md})} 
    
    Defines the mathematical data-generating mechanisms, synthetic epoch timestamps, latent variable representations (e.g., true parameters $\sigma_i$ and $\rho_i$), and analytical metric derivations.
    \item \textbf{Production ORM Schema (\texttt{analysis-resources/production\_schema.md})} 
    
    Formalizes the runtime database constraints, relational foreign-key topologies, and field validators implemented in the production backend (such as strict $1$--$7$ integer Likert bounds on \texttt{app\_emaitemresponse.value\_numeric} enforced via \texttt{MinValueValidator(1)} and \texttt{MaxValueValidator(7)}).
\end{enumerate}
Evaluating simulated trajectories across this dual-schema framework ensures that synthetic validation sets strictly adhere to production persistence constraints while preserving complete visibility into latent affective dynamics.



\section{Simulation System Parameters}

\begin{table}[htbp]
\centering
\small
\begin{tabular}{@{}lp{4.2cm}p{7.8cm}@{}}
\toprule
\textbf{Parameter} & \textbf{Baseline Specification} & \textbf{Methodological Justification} \\
\midrule
Daily EMA Opportunities & 5 prompts/day & Fixed conservative evaluation lower-bound for daily check-ins. \\
Daily Prompt Ceiling & 4 prompts/day & Hard safety cap to bound participant burden. \\
Threshold Quantile & 80th percentile, expanding & Idiographic calibration ensuring out-of-sample thresholding. \\
Intervention Window & 30 minutes & Substantial time to complete the prompt, reducing burden and eliminating pressure for immediate responses.\\
Refractory Cooldown & 60 minutes & Prevents cluster firing during persistent affective episodes. Intervention cooldown.\\
MSSD Rolling Window ($w$) & 3 observations & Rolling statistic capturing localized instability. \\
Burn-in Period & 4 completed days & Emergent requirement ($t \ge 5$) from minimum history requirements. \\
Micro-Randomization ($p$) & 0.50 & Balances intervention power and unconfounded counterfactual estimation. \\
EMA Compliance Baseline & 80\% & Simulation default; distinct from the 75\% preregistered feasibility benchmark and 70\% analytical floor. \\
Deployment Horizon & 35 days (5 weeks) & Primary trial duration (Week 1 run-in followed by 4 intervention weeks). \\
Likert Measurement Bounds & 1 to 7 integer scale & Validator Bounds. \\
\bottomrule
\end{tabular}
\caption{Canonical parameter specification for JITAI simulation experiments.}
\end{table}

\section{Threshold Calibration}

The Mean Square Successive Difference within-person threshold was calibrated to the 80th percentile of each participant's expanding $\widehat{\text{MSSD}}$ vector to optimize the fundamental trade-off between statistical detection power and participant notification burden. In a micro-randomized trial (MRT) featuring up to 5 daily decision points per day and an intervention allocation probability of $p = 0.50$, trial design guidelines require bounding proximal prompt frequency to prevent notification habituation and disengagement while maintaining sufficient variance to detect causal excursion effects.

To empirically calibrate this decision boundary prior to clinical deployment, we conducted a multi-parameter sensitivity simulation crossing a factorial grid of threshold percentiles ($\Theta \in [0.70, 0.90]$ in 0.05 increments) and baseline burn-in durations ($B \in \{0, 1, 2, 3, 5, 7\}$ days) across 150 simulated participants over 10 randomized replications ($N = 1{,}500$ realizations per cell) under varying compliance rates ($R \in \{0.70, 0.80, 0.90\}$). Under the default $80\%$ compliance assumption ($M = 5$ prompts/day) following a 7-day burn-in period, lower percentiles produced excessive intervention dosages of $6.47$ ($\Theta = 0.70$) and $5.43$ ($\Theta = 0.75$) eligible decision points per participant per week. Conversely, conservative thresholds yielded suppressed rates of $3.57$ ($\Theta = 0.85$) and $2.70$ ($\Theta = 0.90$) decision points per week.


\begin{figure}[h]
    \centering
    \includegraphics[width=1.0\textwidth]{threshold_sensitivity_grid_by_resprate.png}
    \caption{Simulated JITAI prompt eligibility across threshold percentiles, burn-in periods, and compliance rates.}
    \label{fig:threshold_sensitivity_grid_by_resprate}
\end{figure}

\begin{figure}[h]
    \centering
    \includegraphics[width=0.8\textwidth]{threshold80_sensitivity_grid.png}
    \caption{Intervention prompt dosage as a function of within-person MSSD threshold percentile and baseline burn-in duration ($R = 80\%$, $n = 1{,}500$).}
    \label{fig:threshold80_sensitivity_grid}
\end{figure}

The 80th-percentile threshold ($\Theta = 0.80$) stabilized empirical dosage at $4.33$ ($\pm 0.05$) eligible decision points per week ($4.62$ at a 3-day burn-in), corresponding to an expected delivery rate of approximately $2.1$ prompts per week post-randomization. This directly satisfies the clinical trial protocol's target dosage of 2--3 delivered prompts per day during high-volatility episodes while respecting the pre-registered upper boundary of 4 prompts per day. Furthermore, factorial sweeps across compliance tiers demonstrated that the response rate acts as an orthogonal scalar: shifting compliance between $70\%$ and $90\%$ modulated the global baseline volume without distorting the underlying relative ranking across percentile contours, establishing $\Theta = 0.80$ as the robust operating cutoff.

\section{Parameter and Metric Recovery}

\subsection{MSSD Tracking Fidelity}

Evaluating an unconstrained cohort ($n = 89$, response rate $R = 0.80$) demonstrates that rolling empirical estimates reliably track latent stationary variance ($r = 0.823$, exceeding the pre-registered construct validity threshold of $r \ge 0.70$). Mean empirical MSSD was $0.953 \pm 0.312$ compared to a latent expectation of $1.155 \pm 0.384$. Across the stationary cohort, $2.6\%$ of all evaluable prompts triggered an eligible intervention point.

\subsection{Autoregressive Parameter Sensitivity}

Latent parameter recovery sweeps demonstrate that recovering innovation variance ($\hat{\sigma}$) is robust across missingness regimes. Conversely, autoregressive persistence ($\hat{\rho}$) degrades non-linearly when EMA completion drops below approximately $70\%$. This establishes $70\%$ as the mathematical breakdown floor for temporal recovery. The trial's $75\%$ pre-registered compliance benchmark incorporates a necessary margin of safety above this floor.

\section{Sensitivity and Robustness Analyses}

\openitem{Sample structure. Need to verify.}

\subsection{Compliance and Burn-in Sensitivity}

A full factorial simulation crossing response rates ($70\% \le R \le 90\%$) with burn-in lengths ($0 \le B \le 7$ days) across 1,500 participant-weeks per cell confirmed that compliance rate and threshold quantile function as orthogonal operational parameters: compliance governs parameter estimation variance, while quantile thresholding directly regulates prompt dosage.

\subsection{Stationarity Misspecification Analysis}

The production engine assumes individual affective volatility is stationary across the 35-day horizon. To evaluate resilience against regime shifts, we evaluated four latent volatility schedules $\sigma_t$ introduced at Day 14[cite: 2]:

\begin{table}[htbp]
\centering
\small
\begin{tabular}{@{}lll@{}}
\toprule
\textbf{Scenario} & \textbf{Mathematical Formulation} & \textbf{Clinical Phenomenon Represented} \\
\midrule
Stationary & $\sigma_t = \sigma_0$ & Unperturbed baseline (control)[cite: 2] \\
Shift Up & $\sigma_t = 2.0\,\sigma_0, \quad t \ge \text{Day } 14$ & Acute exogenous stressor onset[cite: 2] \\
Shift Down & $\sigma_t = 0.5\,\sigma_0, \quad t \ge \text{Day } 14$ & Therapeutic recovery or habituation[cite: 2] \\
Drifting & $\log \sigma_t = \mu_\sigma + 0.97(\log \sigma_{t-1} - \mu_\sigma) + \eta_t$ & Stochastic non-stationary variance drift[cite: 2] \\
\bottomrule
\end{tabular}
\caption{Latent volatility trajectories evaluated in misspecification experiments[cite: 2].}
\end{table}

Performance was benchmarked across 1,500 paired participant realizations (150 participants $\times$ 10 seeds) using common random numbers, comparing the production expanding threshold against a diagnostic 7-day rolling window ($w_{\text{thresh}} = 35$ prompts)[cite: 2]:

\begin{table}[htbp]
\centering
\small
\begin{tabular}{@{}lcccccccc@{}}
\toprule
& \multicolumn{2}{c}{\textbf{Stationary}} & \multicolumn{2}{c}{\textbf{Shift Up}} & \multicolumn{2}{c}{\textbf{Shift Down}} & \multicolumn{2}{c}{\textbf{Drift}} \\
\cmidrule(lr){2-3} \cmidrule(lr){4-5} \cmidrule(lr){6-7} \cmidrule(lr){8-9}
\textbf{Metric} & \textbf{Exp.} & \textbf{Roll.} & \textbf{Exp.} & \textbf{Roll.} & \textbf{Exp.} & \textbf{Roll.} & \textbf{Exp.} & \textbf{Roll.} \\
\midrule
True Positive Rate (TPR)[cite: 2] & 0.51 & 0.52 & 0.70 & 0.61 & 0.26 & 0.37 & 0.46 & 0.48 \\
False Positive Rate (FPR)[cite: 2] & 0.08 & 0.10 & 0.19 & 0.13 & 0.06 & 0.10 & 0.11 & 0.12 \\
Precision[cite: 2] & 0.61 & 0.57 & 0.48 & 0.53 & 0.54 & 0.49 & 0.52 & 0.50 \\
Prompts/Week (Pre-Shift)[cite: 2] & 4.59 & 4.69 & 4.59 & 4.69 & 4.59 & 4.69 & 4.82 & 4.91 \\
Prompts/Week (Post-Shift)[cite: 2] & 4.08 & 4.58 & 9.21 & 6.54 & 1.02 & 3.25 & 4.42 & 4.85 \\
Tracking Ratio (Post-Shift)[cite: 2] & 1.00 & 1.00 & 0.62 & 0.87 & 1.71 & 1.00 & 1.00 & 1.00 \\
Recalibration Latency (Days)[cite: 2] & 0 & 0 & $\infty$ & 5 & $\infty$ & 6 & 0 & 0 \\
\bottomrule
\end{tabular}
\caption{Performance metrics under stationarity misspecification for expanding vs. rolling threshold rules[cite: 2].}
\end{table}

Under an expanding window, the decision engine fails to adapt following discrete variance shifts[cite: 2]. When volatility drops (\textit{Shift Down}), historical data hold the threshold artificially high, causing prompt delivery to collapse to 1.0 prompt/week and reducing true-positive detection to 0.26 without triggering error alerts[cite: 2]. Conversely, when volatility surges (\textit{Shift Up}), prompt dosage increases to 9.2 prompts/week while precision drops to 0.48, generating notification fatigue[cite: 2]. 

The diagnostic 7-day rolling threshold successfully adapts within 5--6 days post-shift, restoring delivery rates at a minor precision loss under stationary conditions (0.61 vs.\ 0.57)[cite: 2]. These findings demonstrate that expanding within-person baselines require automated divergence monitoring to detect regime shifts during clinical field trials[cite: 2].

\end{document}